\chapter{Vitis Platform and GPIO}

\section{Task Description}

Lab 5 uses a standalone MicroBlaze application to read the board switches
through AXI GPIO and write their state to the LED outputs. The software
exercise uses the Vivado hardware export as the input to Vitis.

\section{Project Setup}

The Vitis project contains the firmware source, application settings, and
linker script. Its hardware platform is rebuilt from the Vivado XSA export.
The standalone application starts in \texttt{app.c}. The latest archived
firmware forwards the GPIO input to the dedicated LED controller, as described
in Lab~6.

The archived Lab 5 XSA contains the \texttt{basys3\_top} bitstream and memory
map from Vivado 2025.2.

\section{Lint and Synthesis}

The Lab 5 lint workbook records 36 findings. Four flag unset signal bits,
31 flag unread signal or port bits, and one flags an incomplete case statement
in the debouncer. The unset-bit findings include \texttt{ahb\_hready} at the
board top level. The workbook refers to the Lab 5 revision.

Vivado completed synthesis of \texttt{basys3\_top} with zero errors and zero
critical warnings. The utilization summary reports 2,388 slice LUTs, including
188 used as memory, 2,225 flip-flops, and 32 block RAM tiles.
Table~\ref{tab:resources} gives the archived counts for each synthesis run.
